The eikonal equation imposes an immediate restriction on the zero set of a regular
assignment: its gradient cannot vanish there.  This separates regular zero geometry
from isolated zeros, crossings, and the parameter-dependent singular phenomena
reserved for later work.

\subsection{Regular zero loci}

For an assignment function \(a\), write
\[
 Z(a)=a^{-1}(0).
\]

\begin{theorem}[Regular zero locus]\label{thm:regular-zero-locus}
Let \(\mathfrak E=(M,g,a;\mu,\lambda)\) be a regular AES.
\begin{enumerate}
\item If \(M\) has no boundary, then \(Z(a)\) is a smooth embedded
      codimension-one submanifold of \(M\).
\item If \(M\) has boundary, then
      \(Z(a)\cap\operatorname{int}M\) is a smooth embedded
      codimension-one submanifold of \(\operatorname{int}M\).
\end{enumerate}
In particular, for a surface the indicated zero locus is a disjoint union of
smooth one-dimensional components.  No assertion is made here about finiteness,
properness, compactness, or behavior at \(\partial M\).
\end{theorem}

\begin{proof}
At every \(p\in Z(a)\), the defining eikonal identity gives
\[
 |\nabla a(p)|_g^2=\mu^2.
\]
Because \(\mu\ne0\), one has \(da_p\ne0\).  Thus \(0\) is a regular value of
\(a\) when \(M\) is boundaryless, and of
\(a|_{\operatorname{int}M}\) for interior zeros when a boundary is present.
The regular-value theorem gives the stated embedded submanifolds.
\end{proof}

\begin{corollary}[Local rigidity of regular zeros]\label{cor:regular-zero-rigidity}
Within the interior of a regular two-dimensional AES, the zero locus has no
isolated point, crossing, branch point, or endpoint.  None of these configurations
can occur as the germ of \(Z(a)\) at an interior point.
\end{corollary}

\begin{proof}
By Theorem~\ref{thm:regular-zero-locus}, a neighborhood of every interior zero
admits smooth coordinates in which \(a^{-1}(0)\) is a coordinate line.  Such a
germ has none of the listed configurations.
\end{proof}

This conclusion is local.  It does not classify the number or global topology of
the components.  In the basic hyperbolic model, for example,
\(Z(a)=\{x=0,y>0\}\), but a general regular AES need not have a single zero
component.

Figure~\ref{fig:regular-versus-singular-zero} summarizes the categorical
distinction.  It illustrates the local conclusions already proved; it does not
classify singular zero germs.

\begin{figure}[t]
\centering
\begin{tikzpicture}[
  every node/.style={font=\small},
  grad/.style={-{Latex[length=2mm]},thick,blue!60!black}
]
  \node[font=\bfseries] at (-3.1,1.8) {regular zero};
  \draw[rounded corners=3pt,draw=black!55,fill=blue!3]
    (-5.25,-1.35) rectangle (-0.95,1.35);
  \draw[very thick,red!65!black]
    plot[smooth] coordinates {(-3.9,-1.2) (-3.75,-0.55) (-3.8,0.05)
                              (-3.55,0.65) (-3.65,1.2)};
  \node[red!65!black,anchor=west] at (-3.5,0.85) {$Z(a)$};
  \draw[grad] (-4.7,-0.55) -- (-4.05,-0.55);
  \draw[grad] (-3.45,-0.55) -- (-2.8,-0.55);
  \node[align=center,font=\scriptsize] at (-2.0,-0.75)
    {$|\nabla a|=|\mu|>0$\\smooth zero line};

  \node[font=\bfseries] at (2.8,1.8) {singular isolated zero};
  \fill[orange!3] (2.8,0) circle (1.42);
  \draw[draw=black!55] (2.8,0) circle (1.42);
  \foreach \rad in {0.48,0.88,1.2}
    {\draw[dashed,orange!55!black] (2.8,0) circle (\rad);}
  \fill[red!70!black] (2.8,0) circle (2.2pt);
  \node[anchor=west,red!70!black] at (2.92,0.05) {$p\in S$};
  \node[align=center,font=\scriptsize] at (2.8,-1.75)
    {$a(p)=0$ isolated\\regular axioms not imposed at $p$};
\end{tikzpicture}
\caption{Regular and singular zero geometry.  In a regular AES with
$\mu\ne0$, the assignment crosses zero with nonvanishing gradient, so the
interior zero germ is a smooth curve.  The radial disc model instead has an
isolated zero at its declared singular point, where the assignment is not
$C^1$ and the regular structure has no local extension.}
\label{fig:regular-versus-singular-zero}
\end{figure}

\subsection{Singular arithmetic expression spaces}

The preceding theorem makes a singular framework unavoidable if one wishes to
admit isolated or branching zeros.

\begin{definition}[Singular arithmetic expression space]\label{def:singular-aes}
A \emph{singular arithmetic expression space} is a tuple
\[
 \mathfrak E_{\mathrm{sing}}=(M,S,g,a;\mu,\lambda)
\]
with the following data:
\begin{enumerate}
\item \(M\) is a nonempty oriented smooth surface, possibly with boundary, and
      \(S\subset M\) is a closed nowhere-dense singular set;
\item \(g\) is a smooth Riemannian metric on \(M\setminus S\);
\item \(a:M\to\mathbb R\) is a specified function whose restriction to
      \(M\setminus S\) is smooth;
\item \(\mu\in\mathbb R\setminus\{0\}\), \(\lambda\in\mathbb R\), and
      \[
      (M\setminus S,g,a|_{M\setminus S};\mu,\lambda)
      \]
      is a regular AES.
\item Every \(p\in S\) is locally essential: there is no neighborhood \(U\) of
      \(p\) on which \(g|_{U\setminus S}\) and \(a|_{U\setminus S}\) admit
      smooth extensions \(\widetilde g\) and \(\widetilde a\), with
      \(\widetilde g\) nondegenerate, such that
      \((U,\widetilde g,\widetilde a;\mu,\lambda)\) is a regular AES.
\end{enumerate}
The metric or assignment may possess weaker extensions across \(S\); their precise
regularity and local behavior must be checked in each example.
\end{definition}

The set \(S\) need not consist entirely of zeros.  Nowhere density keeps the regular
locus nonempty and dense, while local essentiality prevents an arbitrary enlargement
of \(S\) by points at which the regular structure extends.  For a singular AES define
\begin{equation}\label{eq:regular-singular-zero-parts}
 Z_{\mathrm{reg}}(a)=Z(a)\cap(M\setminus S),
 \qquad
 Z_{\mathrm{sing}}(a)=Z(a)\cap S.
\end{equation}
Every point of \(Z_{\mathrm{reg}}(a)\) is a regular zero by
Theorem~\ref{thm:regular-zero-locus}; a point of
\(Z_{\mathrm{sing}}(a)\) is called singular because the regular structure is not
asserted there, even when some of the underlying data happen to extend smoothly.

\subsection{A verified isolated-zero singular model}

Let
\[
 \mathbb D=\{(\xi,\eta)\in\mathbb R^2:\xi^2+\eta^2<1\},
 \qquad
 \rho=\sqrt{\xi^2+\eta^2},
\]
and fix \(\mu,\lambda>0\).  Equip the disc with
\begin{equation}\label{eq:isolated-zero-disc-metric}
 g_{\mathbb D}
 =\frac4{\lambda^2(1-\rho^2)^2}
  \bigl(d\xi^2+d\eta^2\bigr)
\end{equation}
and define
\begin{equation}\label{eq:isolated-zero-disc-assignment}
 a_{\mathbb D}(\xi,\eta)
 =
 \begin{cases}
  \displaystyle\frac{2\mu}{\lambda}
  \frac{\rho}{1-\rho^2},&(\xi,\eta)\ne(0,0),\\[7pt]
  0,&(\xi,\eta)=(0,0).
 \end{cases}
\end{equation}

\begin{proposition}[Classification of the isolated-zero disc]\label{prop:isolated-zero-singular-model}
The tuple
\[
 (\mathbb D,\{0\},g_{\mathbb D},a_{\mathbb D};\mu,\lambda)
\]
is a singular AES with the following properties.
\begin{enumerate}
\item The metric \(g_{\mathbb D}\) extends smoothly and
      nondegenerately across the center.
\item The assignment \(a_{\mathbb D}\) is smooth on
      \(\mathbb D\setminus\{0\}\) but is not \(C^1\) at the center.
\item The punctured disc is a regular AES:
      \[
      |\nabla a_{\mathbb D}|_{g_{\mathbb D}}^2
      =\mu^2+\lambda^2a_{\mathbb D}^2.
      \]
\item Its zero decomposition is
      \[
      Z_{\mathrm{reg}}(a_{\mathbb D})=\varnothing,
      \qquad
      Z_{\mathrm{sing}}(a_{\mathbb D})=\{0\}.
      \]
\end{enumerate}
\end{proposition}

\begin{proof}
Formula \eqref{eq:isolated-zero-disc-metric} is written in Cartesian
coordinates.  Its conformal factor is a positive smooth function on all of
\(\mathbb D\), proving the first claim.

Away from the origin, \(\rho\) and hence \(a_{\mathbb D}\) are smooth.  On the
\(\xi\)-axis, however,
\[
 \frac{a_{\mathbb D}(t,0)-a_{\mathbb D}(0,0)}{t}
 =\frac{2\mu}{\lambda}
  \frac{|t|}{t(1-t^2)}.
\]
The one-sided limits are \(2\mu/\lambda\) and
\(-2\mu/\lambda\).  Thus even the partial derivative in the \(\xi\)-direction
does not exist at the center, so the assignment is not \(C^1\) there.  Moreover,
the punctured assignment tends to zero at the center, so its continuous extension
is unique.  The same calculation therefore rules out every smooth local extension
agreeing with the punctured assignment and proves that the declared singular point
is locally essential.

On the punctured disc the assignment is radial, and
\[
 \frac{da_{\mathbb D}}{d\rho}
 =\frac{2\mu}{\lambda}
  \frac{1+\rho^2}{(1-\rho^2)^2}.
\]
The inverse radial metric coefficient is
\[
 g_{\mathbb D}^{\rho\rho}
 =\frac{\lambda^2(1-\rho^2)^2}{4}.
\]
Consequently,
\begin{align*}
 |\nabla a_{\mathbb D}|_{g_{\mathbb D}}^2
 &=g_{\mathbb D}^{\rho\rho}
   \left(\frac{da_{\mathbb D}}{d\rho}\right)^2\\
 &=\mu^2\frac{(1+\rho^2)^2}{(1-\rho^2)^2}\\
 &=\mu^2+
   \lambda^2
   \left(\frac{2\mu}{\lambda}
         \frac{\rho}{1-\rho^2}\right)^2,
\end{align*}
which is the regular-AES equation on
\(\mathbb D\setminus\{0\}\).  Finally,
\(a_{\mathbb D}>0\) for \(0<\rho<1\), while its declared value at the
center is zero.  The claimed zero decomposition follows.
\end{proof}

If \(s=(2/\lambda)\operatorname{arctanh}\rho\) denotes hyperbolic distance
from the center, then
\[
 a_{\mathbb D}=\frac{\mu}{\lambda}\sinh(\lambda s).
\]
This identity describes radial propagation on the punctured regular part; it
does not repair the failure of differentiability at the center.  The metric is
regular there, while the assignment is singular.

\subsection{Smooth parameter families and their total zero sets}

Let \(I\subset\mathbb R\) be an open interval, let \(M\) be a smooth
\(m\)-manifold without boundary, and let
\[
 A:M\times I\longrightarrow\mathbb R
\]
be smooth.  Write \(a_t(p)=A(p,t)\) and
\[
 \mathcal Z=A^{-1}(0)
 =\{(p,t)\in M\times I:a_t(p)=0\}.
\]

\begin{proposition}[Regular total zero set]\label{prop:regular-total-zero-set}
Assume that
\begin{equation}\label{eq:spatial-regular-family-hypothesis}
 d_pa_t\ne0
 \qquad\text{whenever }a_t(p)=0.
\end{equation}
Then \(\mathcal Z\) is a smooth embedded codimension-one submanifold of
\(M\times I\), and the restriction
\begin{equation}\label{eq:zero-family-projection}
 \pi:\mathcal Z\longrightarrow I,
 \qquad
 \pi(p,t)=t,
\end{equation}
is a submersion.  In particular, if \(\dim M=2\), then
\(\mathcal Z\) is a smooth surface.
\end{proposition}

\begin{proof}
At a point \((p,t)\in\mathcal Z\), the total differential is
\[
 dA_{(p,t)}(v,\tau)
 =d_pa_t(v)+d_tA_{(p,t)}(\tau).
\]
Hypothesis \eqref{eq:spatial-regular-family-hypothesis} implies that
\(dA_{(p,t)}\ne0\), so \(0\) is a regular value of \(A\).  The regular-value
theorem proves that \(\mathcal Z\) is a codimension-one embedded submanifold.

It remains to prove the stronger assertion about the projection.  Given an
arbitrary \(\tau\in T_tI\), put
\(c=d_tA_{(p,t)}(\tau)\), and choose \(w\in T_pM\) with
\(d_pa_t(w)=1\), which is possible because \(d_pa_t\ne0\).  Set
\[
 v=-c\,w.
\]
Then \(dA_{(p,t)}(v,\tau)=0\), so
\((v,\tau)\in T_{(p,t)}\mathcal Z\), while
\(d\pi_{(p,t)}(v,\tau)=\tau\).  Hence \(d\pi\) is surjective at every
point of \(\mathcal Z\), proving that \(\pi\) is a submersion.
\end{proof}

\begin{remark}[Properness warning]\label{rem:zero-family-properness}
Proposition~\ref{prop:regular-total-zero-set} gives pointwise product charts
for the projection, but it does not make \(\pi:\mathcal Z\to I\) a globally
trivial family, or even a fiber bundle over an entire parameter neighborhood
when noncompact fibers are present.  A standard sufficient hypothesis for local
triviality as a smooth fiber bundle is that the projection onto the relevant
parameter interval be a proper surjective submersion (with the usual additional
care if boundaries are introduced).  Without properness, components may escape
to or arrive from infinity.  Accordingly, \(\mathcal Z\) is called the
\emph{total zero set} here, not a zero tube.
\end{remark}

No multi-zero model is included in this foundational chapter.  Constructing and
classifying such models requires a complete audit of the domain, metric,
assignment, singular set, flow equation, and exact zero topology; that problem
belongs to the singular-zero theory developed after Paper~I.
